% Standalone reference note for the hybrid-local-solver project.
%
% This document contains only canonical definitions, equivalent formulations,
% and source-backed standard properties. It deliberately contains no
% algorithm, new theorem, experiment, or open conjecture.

\documentclass[11pt]{article}

\usepackage[margin=1in]{geometry}
\usepackage{amsmath,amssymb,amsthm,mathtools,bm}
\usepackage{algorithm}
\usepackage{algorithmic}
\usepackage{booktabs,tabularx,multirow,array,longtable}
\usepackage{enumitem}
\usepackage{microtype}
\usepackage[numbers,sort&compress]{natbib}
\usepackage{xcolor}
\usepackage[colorlinks=true,citecolor=blue,linkcolor=black,urlcolor=blue]{hyperref}
\usepackage[capitalize,nameinlink,noabbrev]{cleveref}

\newtheorem{theorem}{Theorem}[section]
\newtheorem{lemma}[theorem]{Lemma}
\newtheorem{proposition}[theorem]{Proposition}
\newtheorem{corollary}[theorem]{Corollary}
\newtheorem{conjecture}[theorem]{Conjecture}
\theoremstyle{definition}
\newtheorem{definition}[theorem]{Definition}
\newtheorem{assumption}[theorem]{Assumption}
\newtheorem{problem}[theorem]{Problem}
\newtheorem{observation}[theorem]{Observation}
\theoremstyle{remark}
\newtheorem{remark}[theorem]{Remark}
\theoremstyle{plain}

% Canonical mathematical notation for the active manuscript.
%
% This file preserves the recurring notation style used in the author's
% NeurIPS 2024 and 2025 papers while excluding unrelated tensor and
% machine-learning boilerplate from those archived command collections.
%
% Prefix convention:
%   \v...  bold vectors
%   \m...  bold matrices
%   \g...  calligraphic sets, graphs, and families
%   \s...  blackboard-bold spaces and number systems

\newcommand{\method}{\textsc{HybridLocalSolver}}

% Font and scalar shorthands retained from the archived manuscripts.
\newcommand{\mc}{\mathcal}
\newcommand{\R}{\mathbb{R}}
\newcommand{\E}{\mathbb{E}}
\newcommand{\G}{\mathbb{G}}
\newcommand{\N}{\mathcal{N}}
% Accuracy namespaces are semantic: do not add a bare \eps alias.
% Degree-normalized residual tolerance of classical APPR is kept distinct from
% the objective-gap and proximal fixed-point tolerances of the RPPR sections.
\newcommand{\epsappr}{\varepsilon_{\mathrm{appr}}}
\newcommand{\epsobj}{\varepsilon_{\mathrm{obj}}}
\newcommand{\epspg}{\varepsilon_{\mathrm{pg}}}
\newcommand{\epsppr}{\varepsilon_{\mathrm{ppr}}}
\newcommand{\epsin}{\varepsilon_{\mathrm{in}}}
\newcommand{\epskkt}{\varepsilon_{\mathrm{kkt}}}
\newcommand{\epsburn}{\varepsilon_{\mathrm{burn}}}
\newcommand{\epssol}{\varepsilon_{\mathrm{sol}}}
\newcommand{\epspprinst}[1]{\varepsilon_{\mathrm{ppr},#1}}
\newcommand{\trans}{\mathsf{T}}
\newcommand{\grad}{\nabla}

% Delimiters and generic helpers.
% Norms and inner products use fixed-size delimiters by default.
% Use an optional size (e.g. \norm[\big]{...}) or * for automatic sizing.
\DeclarePairedDelimiter{\norm}{\lVert}{\rVert}
\newcommand{\abs}[1]{\left\lvert #1 \right\rvert}
\DeclarePairedDelimiterX{\inner}[2]{\langle}{\rangle}{#1, #2}
\newcommand{\ip}[2]{\inner{#1}{#2}}
\newcommand{\card}[1]{\left\lvert #1 \right\rvert}
\newcommand{\set}[1]{\left\{#1\right\}}
\newcommand{\ceil}[1]{\left\lceil #1 \right\rceil}
\newcommand{\floor}[1]{\left\lfloor #1 \right\rfloor}
\newcommand{\comp}[1]{\overline{#1}}
\newcommand{\conj}[1]{\overline{#1}}
\newcommand{\mean}[1]{\overline{#1}}
\newcommand{\bigO}{\mathcal{O}}
\newcommand{\softO}{\widetilde{\mathcal{O}}}
\newcommand{\softOmega}{\widetilde{\Omega}}
\newcommand{\pos}[1]{\left[#1\right]_{+}}

% Bold lowercase vectors.
\newcommand{\va}{\bm{a}}
\newcommand{\vb}{\bm{b}}
\newcommand{\vc}{\bm{c}}
\newcommand{\vd}{\bm{d}}
\newcommand{\ve}{\bm{e}}
\newcommand{\vf}{\bm{f}}
\newcommand{\vg}{\bm{g}}
\newcommand{\vh}{\bm{h}}
\newcommand{\vi}{\bm{i}}
\newcommand{\vj}{\bm{j}}
\newcommand{\vk}{\bm{k}}
\newcommand{\vl}{\bm{l}}
\newcommand{\vm}{\bm{m}}
\newcommand{\vn}{\bm{n}}
\newcommand{\vo}{\bm{o}}
\newcommand{\vp}{\bm{p}}
\newcommand{\vq}{\bm{q}}
\newcommand{\vr}{\bm{r}}
\newcommand{\vs}{\bm{s}}
\newcommand{\vt}{\bm{t}}
\newcommand{\vu}{\bm{u}}
\newcommand{\vv}{\bm{v}}
\newcommand{\vw}{\bm{w}}
\newcommand{\vx}{\bm{x}}
\newcommand{\vy}{\bm{y}}
\newcommand{\vz}{\bm{z}}

% Bold Greek vectors and special vectors.
\newcommand{\valpha}{\bm{\alpha}}
\newcommand{\vbeta}{\bm{\beta}}
\newcommand{\vdelta}{\bm{\delta}}
\newcommand{\vepsilon}{\bm{\epsilon}}
\newcommand{\veta}{\bm{\eta}}
\newcommand{\vgamma}{\bm{\gamma}}
\newcommand{\vlambda}{\bm{\lambda}}
\newcommand{\vmu}{\bm{\mu}}
\newcommand{\vomega}{\bm{\omega}}
\newcommand{\vphi}{\bm{\phi}}
\newcommand{\vpi}{\bm{\pi}}
\newcommand{\vpsi}{\bm{\psi}}
\newcommand{\vrho}{\bm{\rho}}
\newcommand{\vsigma}{\bm{\sigma}}
\newcommand{\vtheta}{\bm{\theta}}
\newcommand{\vxi}{\bm{\xi}}
\newcommand{\vell}{\bm{\ell}}
\newcommand{\vDelta}{\bm{\Delta}}
\newcommand{\vGamma}{\bm{\Gamma}}
\newcommand{\vLambda}{\bm{\Lambda}}
\newcommand{\vPhi}{\bm{\Phi}}
\newcommand{\vSigma}{\bm{\Sigma}}
\newcommand{\vzero}{\bm{0}}
\newcommand{\vone}{\bm{1}}
\newcommand{\one}{\mathbf{1}}
\newcommand{\zero}{\vzero}
\newcommand{\eunit}[1]{\bm{e}_{#1}}
\newcommand{\vDeltax}{\bm{\Delta}\bm{x}}

% Bold uppercase matrices.
\newcommand{\mA}{\bm{A}}
\newcommand{\mB}{\bm{B}}
\newcommand{\mC}{\bm{C}}
\newcommand{\mD}{\bm{D}}
\newcommand{\mE}{\bm{E}}
\newcommand{\mF}{\bm{F}}
\newcommand{\mG}{\bm{G}}
\newcommand{\mH}{\bm{H}}
\newcommand{\mI}{\bm{I}}
\newcommand{\mJ}{\bm{J}}
\newcommand{\mK}{\bm{K}}
\newcommand{\mL}{\bm{L}}
\newcommand{\mM}{\bm{M}}
\newcommand{\mN}{\bm{N}}
\newcommand{\mO}{\bm{O}}
\newcommand{\mP}{\bm{P}}
\newcommand{\mQ}{\bm{Q}}
\newcommand{\mR}{\bm{R}}
\newcommand{\mS}{\bm{S}}
\newcommand{\mT}{\bm{T}}
\newcommand{\mU}{\bm{U}}
\newcommand{\mV}{\bm{V}}
\newcommand{\mW}{\bm{W}}
\newcommand{\mX}{\bm{X}}
\newcommand{\mY}{\bm{Y}}
\newcommand{\mZ}{\bm{Z}}

% Bold Greek matrices retained from the graph papers.
\newcommand{\mBeta}{\bm{\beta}}
\newcommand{\mDelta}{\bm{\Delta}}
\newcommand{\mGamma}{\bm{\Gamma}}
\newcommand{\mLambda}{\bm{\Lambda}}
\newcommand{\mOmega}{\bm{\Omega}}
\newcommand{\mPhi}{\bm{\Phi}}
\newcommand{\mPi}{\bm{\Pi}}
\newcommand{\mSigma}{\bm{\Sigma}}
\newcommand{\mTheta}{\bm{\Theta}}

% Calligraphic symbols.
\newcommand{\gA}{\mathcal{A}}
\newcommand{\gB}{\mathcal{B}}
\newcommand{\gC}{\mathcal{C}}
\newcommand{\gD}{\mathcal{D}}
\newcommand{\gE}{\mathcal{E}}
\newcommand{\gF}{\mathcal{F}}
\newcommand{\gG}{\mathcal{G}}
\newcommand{\gH}{\mathcal{H}}
\newcommand{\gI}{\mathcal{I}}
\newcommand{\gJ}{\mathcal{J}}
\newcommand{\gK}{\mathcal{K}}
\newcommand{\gL}{\mathcal{L}}
\newcommand{\gM}{\mathcal{M}}
\newcommand{\gN}{\mathcal{N}}
\newcommand{\gO}{\mathcal{O}}
\newcommand{\gP}{\mathcal{P}}
\newcommand{\gQ}{\mathcal{Q}}
\newcommand{\gR}{\mathcal{R}}
\newcommand{\gS}{\mathcal{S}}
\newcommand{\gT}{\mathcal{T}}
\newcommand{\gU}{\mathcal{U}}
\newcommand{\gV}{\mathcal{V}}
\newcommand{\gW}{\mathcal{W}}
\newcommand{\gX}{\mathcal{X}}
\newcommand{\gY}{\mathcal{Y}}
\newcommand{\gZ}{\mathcal{Z}}

% Blackboard-bold symbols.
\newcommand{\sA}{\mathbb{A}}
\newcommand{\sB}{\mathbb{B}}
\newcommand{\sC}{\mathbb{C}}
\newcommand{\sD}{\mathbb{D}}
\newcommand{\sE}{\mathbb{E}}
\newcommand{\sF}{\mathbb{F}}
\newcommand{\sG}{\mathbb{G}}
\newcommand{\sH}{\mathbb{H}}
\newcommand{\sI}{\mathbb{I}}
\newcommand{\sJ}{\mathbb{J}}
\newcommand{\sK}{\mathbb{K}}
\newcommand{\sL}{\mathbb{L}}
\newcommand{\sM}{\mathbb{M}}
\newcommand{\sN}{\mathbb{N}}
\newcommand{\sO}{\mathbb{O}}
\newcommand{\sP}{\mathbb{P}}
\newcommand{\sQ}{\mathbb{Q}}
\newcommand{\sR}{\mathbb{R}}
\newcommand{\sS}{\mathbb{S}}
\newcommand{\sT}{\mathbb{T}}
\newcommand{\sU}{\mathbb{U}}
\newcommand{\sV}{\mathbb{V}}
\newcommand{\sW}{\mathbb{W}}
\newcommand{\sX}{\mathbb{X}}
\newcommand{\sY}{\mathbb{Y}}
\newcommand{\sZ}{\mathbb{Z}}

% Frequently used operators.
\DeclareMathOperator{\Cov}{Cov}
\DeclareMathOperator{\Var}{Var}
\DeclareMathOperator{\diag}{diag}
\DeclareMathOperator{\nei}{Nei}
\DeclareMathOperator{\nnz}{nnz}
\DeclareMathOperator{\pr}{pr}
\DeclareMathOperator{\prox}{prox}
\DeclareMathOperator{\push}{push}
\DeclareMathOperator{\res}{res}
\DeclareMathOperator{\rel}{Rel}
\DeclareMathOperator{\relu}{ReLU}
\DeclareMathOperator{\sign}{sign}
\DeclareMathOperator{\supp}{supp}
\DeclareMathOperator{\tr}{tr}
\DeclareMathOperator{\Tr}{Tr}
\DeclareMathOperator{\vol}{vol}
\DeclareMathOperator{\work}{work}
\DeclareMathOperator{\Work}{Work}
\DeclareMathOperator*{\argmax}{arg\,max}
\DeclareMathOperator*{\argmin}{arg\,min}

% Project-wide names and claim-status labels used by standalone research notes.
% Mathematical notation belongs in math_commands.tex.

% Algorithms and solver families.
\newcommand{\AESP}{\textsc{AESP}}
\newcommand{\APPR}{\textsc{APPR}}
\newcommand{\ASPR}{\textsc{ASPR}}
\newcommand{\FISTA}{\textsc{FISTA}}
\newcommand{\ISTA}{\textsc{ISTA}}
\newcommand{\LOCAPPR}{\textsc{LocAPPR}}
\newcommand{\LOCGD}{\textsc{LocGD}}
\newcommand{\LOCSOR}{\textsc{LocSOR}}
\newcommand{\LOCCH}{\textsc{LocCH}}
\newcommand{\LOCHB}{\textsc{LocHB}}
\newcommand{\LocProxGS}{\textsc{LocProxGS}}
\newcommand{\Hybrid}{\textsc{AESP--LocSOR}}

% Claim classifications required by the repository research-note protocol.
\newcommand{\statussource}{\textbf{Source result}}
\newcommand{\statusproved}{\textbf{Proved here}}
\newcommand{\statusconditional}{\textbf{Conditional}}
\newcommand{\statusconjecture}{\textbf{Open / conjectural}}
\newcommand{\statusopen}{\textbf{Open}}
\newcommand{\statusempirical}{\textbf{Empirical}}
\newcommand{\statusobservation}{\textbf{Empirical observation}}
\newcommand{\statusrefuted}{\textbf{Refuted route}}

% Compatibility names used by the older complete-note presentation layer.
\newcommand{\Status}[2]{\textbf{#1:} #2}
\newcommand{\SourceStatus}{\textnormal{\textsc{Source result}}}
\newcommand{\ProvedStatus}{\textnormal{\textsc{Proved here}}}
\newcommand{\ConditionalStatus}{\textnormal{\textsc{Conditional}}}
\newcommand{\ComputedStatus}{\textnormal{\textsc{Computed}}}
\newcommand{\RefutedStatus}{\textnormal{\textsc{Refuted route}}}
\newcommand{\OpenStatus}{\textnormal{\textsc{Open}}}


\title{Problem Definitions and Standard Properties\\
for Local PageRank and Regularized PageRank}
\author{Baojian Zhou\\
\small Working reference note for \texttt{hybrid-local-solver}}
\date{August 2026}

\begin{document}
\maketitle

\begin{abstract}
This note fixes a compact reference for the two optimization problems used in
the project: personalized PageRank (PPR) and regularized personalized PageRank
(RPPR).  It records the exact graph, seed, parameter, output, and local-access
models; translates the lazy symmetric system into lazy fixed-point, rescaled,
non-lazy mass, and non-lazy degree coordinates; and collects standard
spectral, positivity, normalization, optimality, sparsity, bias, and residual
properties.  Every substantive property is either a cited source result or an
elementary consequence of the displayed definitions.  No algorithm,
complexity claim, experiment, or new research statement is included, and the
repository-wide residual decision remains open.
\end{abstract}

\section{Scope and status}
\label{sec:problem-definitions-scope}

\paragraph{Claim status.}
This is a \statussource{} reference document.  The shared block below is the
canonical manuscript definition; this note does not replace or redefine it.
Statements labeled ``standard consequence'' are included only to make the
reference self-contained.  They are direct linear-algebra or convex-analysis
consequences of the cited source facts, not claims of novelty.

\paragraph{Included objects.}
The note covers the PPR linear system and fixed point, the RPPR convex
surrogate, semantic output accuracy, a sufficient residual-to-error
implication, and the adjacency-list input and work unit.  It excludes every
specific solver, activation rule, momentum scheme, convergence rate,
complexity target, experiment, and open conjecture.

\subsection{Common source-aligned PageRank and RPPR model}
\label{subsec:shared-source-aligned-problem}

\textbf{Primary notation.}
All documents in this manuscript workspace use the following reference
language. Bold lowercase letters denote column vectors, and bold uppercase
letters denote matrices. The displayed domains determine the dimensions.
A note may impose a stronger assumption after this block, but it must not
redefine these objects. Let $\mathcal{G}=(\mathcal{V},\mathcal{E})$ be a
finite, simple, undirected, unweighted graph without isolated vertices. Let
$\mathcal{V}=[n]:=\{1,\ldots,n\}$, and let
$\bm{A}\in\{0,1\}^{n\times n}$ be its symmetric adjacency matrix. For
$i\in\mathcal{V}$, let
$\mathcal{N}(i):=\{j\in\mathcal{V}:\{i,j\}\in\mathcal{E}\}$, define
$d_i:=|\mathcal{N}(i)|$, and put
$\bm{D}:=\diag(d_1,\ldots,d_n)$.

For $\mathcal{S}\subseteq\mathcal{V}$, write
$\mathcal{N}(\mathcal{S}):=\bigcup_{i\in\mathcal{S}}\mathcal{N}(i)$ and
$\partial\mathcal{S}:=\mathcal{N}(\mathcal{S})\setminus\mathcal{S}$. Define
$\operatorname{Ext}(\mathcal{S}):=\mathcal{V}\setminus
(\mathcal{S}\cup\partial\mathcal{S})$ and
$\vol(\mathcal{S}):=\sum_{i\in\mathcal{S}}d_i$. The support of a vector
$\bm{x}$ is $\supp(\bm{x}):=\{i\in\mathcal{V}:x_i\neq0\}$.
Matrix inequalities use the Loewner order, and all unqualified vector
inequalities are coordinatewise.

The seed is a column distribution $\bm{s}\in\R_+^n$ satisfying
$\inner{\one}{\bm{s}}=1$.
The single-seed specialization is always written $\bm{s}=\bm{e}_v$, where
$v\in\mathcal{V}$ is the seed vertex. Thus, $\bm{s}$ is never used as a vertex
label. For $\alpha\in(0,1]$, define
\begin{equation}
\begin{aligned}
    \bm{\mathcal{L}}
        &:=\bm{I}-\bm{D}^{-1/2}\bm{A}\bm{D}^{-1/2},\\
    \bm{Q}
        &:=\alpha\bm{I}+\frac{1-\alpha}{2}\bm{\mathcal{L}}
          =\frac{1+\alpha}{2}\bm{I}
           -\frac{1-\alpha}{2}
            \bm{D}^{-1/2}\bm{A}\bm{D}^{-1/2},\\
    \bm{b}
        &:=\alpha\bm{D}^{-1/2}\bm{s}.
\end{aligned}
    \label{eq:shared-pagerank-matrices}
\end{equation}
Since $\bm{0}\preceq\bm{\mathcal{L}}\preceq2\bm{I}$,
$\alpha\bm{I}\preceq\bm{Q}\preceq\bm{I}$. The shared smooth PageRank
quadratic is
\begin{equation}
    f(\bm{x})
    :=\frac12\inner{\bm{x}}{\bm{Q}\bm{x}}
      -\inner{\bm{b}}{\bm{x}},
    \qquad
    \nabla f(\bm{x})=\bm{Q}\bm{x}-\bm{b}.
    \label{eq:shared-pagerank-objective}
\end{equation}
Its unique unregularized minimizer and associated lazy personalized PageRank
vector are
\begin{equation}
    \bm{x}_0^*:=\bm{Q}^{-1}\bm{b},
    \qquad
    \bm{\pi}:=\bm{D}^{1/2}\bm{x}_0^*.
    \label{eq:shared-ppr-solution}
\end{equation}

For a regularization parameter $\rho>0$, reserve $g_\rho$ for the nonsmooth
weighted $\ell_1$ term and define
\begin{equation}
    g_\rho(\bm{x}):=\alpha\rho\|\bm{D}^{1/2}\bm{x}\|_1,
    \qquad
    F_\rho(\bm{x}):=f(\bm{x})+g_\rho(\bm{x}),
    \qquad
    \bm{x}_\rho^*
    :=\argmin_{\bm{x}\in\R^n}F_\rho(\bm{x}).
    \label{eq:shared-rppr-objective}
\end{equation}
When $\rho = 0$, it is the unregularized point $\bm{x}_0^*$. Write
$\mathcal{S}_\rho^* := \supp(\bm{x}_\rho^*)$ and
$\mathcal{I}_\rho^* := [n]\setminus\mathcal{S}_\rho^*$.
The source records $\bm{x}_\rho^*\geq\bm{0}$,
$\vol(\mathcal{S}_\rho^*)\leq1/\rho$, and the coordinatewise conditions
\begin{equation}
    \nabla_i f(\bm{x}_\rho^*)
    \in
    \begin{cases}
        \{-\alpha\rho\sqrt{d_i}\}, & x_i^*(\rho)>0,\\
        [-\alpha\rho\sqrt{d_i},0], & x_i^*(\rho)=0.
    \end{cases}
    \label{eq:shared-rppr-kkt}
\end{equation}

The common computational unit is one repeated adjacency-list scan: scanning
coordinate $i$ costs $d_i$, and scanning a set $\mathcal{S}$ costs
$\vol(\mathcal{S})$.
Iteration count and wall-clock time do not replace this degree-weighted work
measure.

Finally, accuracy symbols remain distinct. The APPR threshold is
$\epsappr$; $\epsobj$ denotes an objective-gap target; $\epspg$ denotes the
source experiment's proximal fixed-point tolerance; and $\epsppr$ may denote
a document-scoped degree-normalized PPR error target.
No equality or conversion among these quantities is assumed. In particular,
this shared model does not resolve the repository-wide residual decision.


\section{Exact computational problems}
\label{sec:problem-definitions-problems}

\begin{problem}[Personalized PageRank]
\label{prob:problem-definitions-ppr}
The input is the graph $\mathcal{G}$ and sparse seed distribution $\bm{s}$
from the shared model, together with $\alpha\in(0,1]$. The exact target is
$\bm{x}^0=\bm{Q}^{-1}\bm{b}$, equivalently
$\bm{\pi}=\bm{D}^{1/2}\bm{x}^0$. A sparse approximate output is a list
representing $\widehat{\bm{x}}\in\R^n$; unlisted coordinates are zero and
$\widehat{\bm{\pi}}:=\bm{D}^{1/2}\widehat{\bm{x}}$. For a
document-scoped tolerance $\epsppr>0$, the semantic accuracy requirement is
\begin{equation}
    \norm{\bm{D}^{-1}
      (\widehat{\bm{\pi}}-\bm{\pi})}_\infty
    =
    \norm{\bm{D}^{-1/2}
      (\widehat{\bm{x}}-\bm{x}^0)}_\infty
    \leq \epsppr.
    \label{eq:problem-definitions-semantic-ppr}
\end{equation}
This is an output property, not an activation rule or an adopted
implementation-wide stopping test.
\end{problem}

\begin{problem}[Regularized personalized PageRank]
\label{prob:problem-definitions-rppr}
For the same $\mathcal{G},\bm{s},\alpha$ and a regularization parameter
$\rho>0$, the exact target is the unique minimizer $\bm{x}^\star(\rho)$ of
$F_\rho$ in \cref{eq:shared-rppr-objective}. Its mass-coordinate
representation is $\bm{\pi}_\rho:=\bm{D}^{1/2}\bm{x}^\star(\rho)$, with
$\mathcal{S}^\star(\rho)=\supp(\bm{x}^\star(\rho))$.
The scale $\rho$ is part of the objective and is not silently identified with
$\epsppr$, an objective-gap tolerance, or an algorithmic KKT diagnostic.
\end{problem}

The input graph is accessed through adjacency lists and the seed through its
sparse nonzero list.  Scanning vertex $i$ costs $d_i$ and scanning
$\mathcal{S}\subseteq\mathcal{V}$ costs $\vol(\mathcal{S})$. This defines
the common local-work unit; it does not prescribe an algorithm or assert a
work bound.

\section{Equivalent PPR formulations}
\label{sec:problem-definitions-equivalences}

Put $\bm{P}:=\bm{A}\bm{D}^{-1}$,
$\bm{P}_{\rm L}:=(\bm{I}+\bm{P})/2$, and
$\alpha_{\rm nl}:=2\alpha/(1+\alpha)$.
Both $\bm{P}$ and $\bm{P}_{\rm L}$ act on column vectors, and both are
column stochastic. Direct multiplication of
$\bm{Q}\bm{x}^0=\bm{b}$ by $\bm{D}^{1/2}$ gives the exact dictionary below.
The lazy forms agree with the PPR formulations used by
\citet{zhou2024iterative} and \citet{huang2025accelerated}; the non-lazy mass
and degree forms are the undirected specializations used by
\citet{chen2023accelerating} and \citet{wei2026simple}.

\begin{center}
\begin{tabularx}{0.97\textwidth}{@{}
    >{\raggedright\arraybackslash}p{0.22\textwidth}
    >{\raggedright\arraybackslash}X@{}}
\toprule
Form & Exact equation and conversion \\
\midrule
Lazy fixed point
& $\bm{\pi}=\alpha\bm{s}
   +(1-\alpha)\bm{P}_{\rm L}\bm{\pi}$. \\
Lazy symmetric
& $\bm{Q}\bm{x}^0=\bm{b}$ and
  $\bm{\pi}=\bm{D}^{1/2}\bm{x}^0$. \\
Rescaled lazy symmetric
& $\bm{Q}_{\rm rs}\bm{x}^0=\bm{b}_{\rm rs}$, where
  $\bm{Q}_{\rm rs}=2\bm{Q}/(1+\alpha)$ and
  $\bm{b}_{\rm rs}=2\bm{b}/(1+\alpha)$. \\
Non-lazy mass
& $[\bm{I}-(1-\alpha_{\rm nl})\bm{P}]\bm{\pi}
   =\alpha_{\rm nl}\bm{s}$, where
  $\alpha_{\rm nl}=2\alpha/(1+\alpha)$ and
  $\alpha=\alpha_{\rm nl}/(2-\alpha_{\rm nl})$. \\
Non-lazy degree
& $[\bm{D}-(1-\alpha_{\rm nl})\bm{A}]\bm{y}
   =\alpha_{\rm nl}\bm{s}$, where
  $\bm{y}=\bm{D}^{-1}\bm{\pi}
   =\bm{D}^{-1/2}\bm{x}^0$. \\
\bottomrule
\end{tabularx}
\end{center}

The rescaled system changes neither $\alpha$ nor its solution.  Removing the
lazy self-loop changes the numerical teleportation parameter to
$\alpha_{\rm nl}$.  Therefore a dependence on ``$\alpha$'' can be compared
across sources only after applying this conversion.

\section{Standard properties of PPR}
\label{sec:problem-definitions-ppr-properties}

The following facts are standard for the normalized Laplacian and the
discounted PageRank system; see the source formulations and spectral bounds
in \citet{fountoulakis2019variational,fountoulakis2026complexity,
zhou2024iterative,huang2025accelerated}.

\begin{proposition}[Spectrum, convexity, and uniqueness; standard]
\label{prop:problem-definitions-spectrum}
The normalized Laplacian and PageRank matrix satisfy
\begin{equation}
    \bm{0}\preceq\bm{\mathcal{L}}\preceq2\bm{I},
    \qquad
    \alpha\bm{I}\preceq\bm{Q}\preceq\bm{I}.
    \label{eq:problem-definitions-spectrum}
\end{equation}
Consequently, $\bm{Q}$ is symmetric positive definite,
$\kappa_2(\bm{Q})\leq1/\alpha$, $f$ is $\alpha$-strongly convex with
$1$-Lipschitz gradient, and $\bm{x}^0$ is unique. Moreover,
\begin{equation}
    f(\bm{x})-f(\bm{x}^0)
    =\frac12\norm{\bm{x}-\bm{x}^0}_{\bm{Q}}^2
    \qquad(\bm{x}\in\R^n).
    \label{eq:problem-definitions-quadratic-gap}
\end{equation}
\end{proposition}

\begin{proposition}[Inverse positivity and probability mass; standard]
\label{prop:problem-definitions-positivity}
The matrix $\bm{Q}$ is a nonsingular symmetric $M$-matrix (a Stieltjes
matrix), so $\bm{Q}^{-1}\geq\bm{0}$ coordinatewise. Hence
$\bm{x}^0\geq\bm{0}$ and $\bm{\pi}\geq\bm{0}$. In addition,
\begin{equation}
    \inner{\one}{\bm{\pi}}=1.
    \label{eq:problem-definitions-unit-mass}
\end{equation}
If $0<\alpha<1$, then $\bm{\pi}$ is strictly positive on every connected
component that intersects $\supp(\bm{s})$. At $\alpha=1$, one simply has
$\bm{\pi}=\bm{s}$.
\end{proposition}

\begin{proof}[Standard verification]
The off-diagonal entries of $\bm{Q}$ are nonpositive, and
\cref{eq:problem-definitions-spectrum} gives positive definiteness; this is
the Stieltjes characterization.  The lazy fixed point and column
stochasticity of $\bm{P}_{\rm L}$ give
$\one^T\bm{\pi}=\alpha+(1-\alpha)\one^T\bm{\pi}$. Strict positivity on a seeded
component follows from the Neumann series for the discounted walk.
\end{proof}

\begin{proposition}[A sufficient degree-normalized certificate; standard]
\label{prop:problem-definitions-residual-certificate}
For every $\widehat{\bm{x}}\in\R^n$,
\begin{equation}
    \norm{\bm{D}^{-1}
      (\widehat{\bm{\pi}}-\bm{\pi})}_\infty
    \leq
    \frac{1}{\alpha}
    \norm{\bm{D}^{-1/2}
      (\bm{Q}\widehat{\bm{x}}-\bm{b})}_\infty.
    \label{eq:problem-definitions-residual-certificate}
\end{equation}
Thus
$\norm{\bm{D}^{-1/2}
  (\bm{Q}\widehat{\bm{x}}-\bm{b})}_\infty\leq\alpha\epsppr$
is sufficient for \cref{eq:problem-definitions-semantic-ppr}.  This is the
scaling used in the source-aligned local certificate of
\citet{zhou2024iterative}; it does not select the sign, strictness, or
evaluation schedule of a repository-wide residual.
\end{proposition}

\begin{proof}[Standard verification]
Let $\bm{B}:=\bm{D}^{-1/2}\bm{Q}\bm{D}^{1/2}$. It is a nonsingular
$M$-matrix, $\bm{B}^{-1}\geq\bm{0}$, and
$\bm{B}\one=\alpha\one$. Therefore
$\norm{\bm{B}^{-1}}_\infty=1/\alpha$. Applying $\bm{B}^{-1}$ to the scaled
residual gives \cref{eq:problem-definitions-residual-certificate}.
\end{proof}

\section{Standard properties of RPPR}
\label{sec:problem-definitions-rppr-properties}

The source variational formulation, KKT system, nonnegative minimizer, and
support-volume result are given in \citet[Section~4 and
Theorems~1--2]{fountoulakis2019variational}; the shared normalization follows
\citet[Section~3]{fountoulakis2026complexity}.

\begin{proposition}[Existence, optimality, and locality; source]
\label{prop:problem-definitions-rppr-source}
For every $\rho>0$, $F_\rho$ is $\alpha$-strongly convex and has one minimizer
$\bm{x}^\star(\rho)$. This minimizer is nonnegative, obeys the coordinatewise
KKT conditions in \cref{eq:shared-rppr-kkt}, and satisfies
\begin{equation}
    \vol(\mathcal{S}^\star(\rho))\leq\frac{1}{\rho}.
    \label{eq:problem-definitions-support-volume}
\end{equation}
In particular,
$F_\rho(\bm{x})-F_\rho(\bm{x}^\star)\geq
\frac{\alpha}{2}\norm{\bm{x}-\bm{x}^\star}_2^2$ for all $\bm{x}$.
\end{proposition}

\begin{proposition}[Regularization bias; standard consequence]
\label{prop:problem-definitions-rppr-bias}
Let $\bm{\pi}_\rho=\bm{D}^{1/2}\bm{x}^\star(\rho)$. Then
\begin{equation}
    \bm{0}
    \leq \bm{D}^{-1}(\bm{\pi}-\bm{\pi}_\rho)
    \leq \rho\one,
    \qquad
    \norm{\bm{D}^{-1}(\bm{\pi}-\bm{\pi}_\rho)}_\infty\leq\rho.
    \label{eq:problem-definitions-rppr-bias}
\end{equation}
Consequently, any approximation $\widehat{\bm{x}}$ satisfying
\begin{equation}
    \norm{\bm{D}^{-1/2}
      (\widehat{\bm{x}}-\bm{x}^\star(\rho))}_\infty\leq\tau
    \label{eq:problem-definitions-rppr-optimization-error}
\end{equation}
obeys the PPR error bound
\begin{equation}
    \norm{\bm{D}^{-1}
      (\bm{D}^{1/2}\widehat{\bm{x}}-\bm{\pi})}_\infty
    \leq\rho+\tau.
    \label{eq:problem-definitions-rppr-total-error}
\end{equation}
\end{proposition}

\begin{proof}[Standard verification]
The KKT conditions provide a vector $\bm{z}\in[0,1]^n$ such that
$\bm{Q}\bm{x}^\star-\bm{b}
=-\alpha\rho\bm{D}^{1/2}\bm{z}$. Since
$\bm{Q}\bm{x}^0=\bm{b}$,
\[
    \bm{x}^0-\bm{x}^\star
    =\alpha\rho\bm{Q}^{-1}\bm{D}^{1/2}\bm{z}.
\]
Inverse positivity and
$\bm{Q}\bm{D}^{1/2}\one=\alpha\bm{D}^{1/2}\one$ give
$\bm{0}\leq\bm{x}^0-\bm{x}^\star
\leq\rho\bm{D}^{1/2}\one$, which is
\cref{eq:problem-definitions-rppr-bias}.  The last bound is the triangle
inequality.
\end{proof}

Thus $\rho$ has two separate, standard consequences: it bounds the exact
degree-normalized bias and it bounds the optimal support volume.  For example,
choosing $\rho=\epsppr/2$ leaves an error budget $\tau=\epsppr/2$ for an
inexact RPPR solve and gives
$\vol(\mathcal{S}^\star)\leq2/\epsppr$. This is only a problem-level
implication; it does not supply an algorithm for attaining $\tau$.

\section{Notation and source boundary}
\label{sec:problem-definitions-source-boundary}

The symbols below are intentionally noninterchangeable.

\begin{center}
\begin{tabularx}{0.97\textwidth}{@{}
    >{\raggedright\arraybackslash}p{0.17\textwidth}
    >{\raggedright\arraybackslash}X@{}}
\toprule
Symbol & Meaning \\
\midrule
$\epsppr$ & Semantic degree-normalized PPR output error in
\cref{eq:problem-definitions-semantic-ppr}. \\
$\epsappr$ & APPR residual activation threshold, defined only in the APPR
algorithm; it is not $\epsppr$. \\
$\rho$ & RPPR regularization and locality scale. \\
$\epsobj$ & Additive RPPR objective-gap target. \\
$\epspg$ & Proximal fixed-point diagnostic used by source experiments. \\
$\tau$ & A local symbol for the RPPR-to-PPR optimization-error allowance in
\cref{eq:problem-definitions-rppr-optimization-error}. \\
\bottomrule
\end{tabularx}
\end{center}

The exact source audit trail is maintained in the literature notes:
Fountoulakis et al.\ (2019), PDF pages 8--18 for the variational objective,
KKT system, nonnegativity, support bound, and relative stopping rule;
Fountoulakis and Mart\'inez-Rubio (2026), PDF page 3, Section 3 for the shared
RPPR normalization; Zhou et al.\ (2024), PDF pages 1--3 and supplemental PDF
page 17 for the lazy PPR systems and residual certificate; Huang et al.\
(2025), PDF pages 1--3 for the lazy PPR and symmetric quadratic; Chen et al.\
(2023), PDF page 3 for the non-lazy mass system; and Wei and Yang (2026), PDF
pages 2--4 and 8--10 for the non-lazy degree-coordinate system and active-set
RPPR properties.  These pointers are recorded in
\texttt{docs/literature/local-solvers.md} and
\texttt{docs/literature/acceleration.md}.

\paragraph{Boundary of the note.}
Nothing here chooses a solver, claims graph-uniform locality, identifies two
accuracy namespaces, or resolves the implementation-wide residual convention.
Those choices require their own definitions, decisions, and proofs.

\bibliographystyle{plainnat}
\bibliography{../../references}

\end{document}
